\documentclass[../../main.tex]{subfiles}
\graphicspath{{\subfix{../../image/}}}

\begin{document}

\section{事件独立性}\label{section:ch02-事件独立性}

\subsection{两个事件的独立性}

本节中我们引进一个新的概念——统计独立性. 先从两个事件的独立性开始, 然后讨论更为一般的场合.

还是从考虑古典概型的一个例子作为出发点.

\begin{example}[有放回摸球]\label{example:ch02-有放回摸球}
一口袋中装有 $a$ 只黑球和 $b$ 只白球, 采用有放回摸球, 求：

（1）在已知第一次摸得黑球的条件下, 第二次摸出黑球的概率；

（2）第二次摸出黑球的概率.

\begin{proof}[解]
以事件 $A$ 表示第一次摸得黑球, 事件 $B$ 表示第二次摸得黑球. 则
\begin{align*}
P(A)=\frac{a}{a+b},\quad P(AB)=\frac{a^2}{(a+b)^2},\quad P(\bar{A}B)=\frac{ba}{(a+b)^2}
\end{align*}
所以
\begin{align*}
P(B\mid A)=\frac{P(AB)}{P(A)}=\frac{a}{a+b}
\end{align*}
而
\begin{align*}
P(B)=P(AB)+P(\bar{A}B)=\frac{a^2}{(a+b)^2}+\frac{ba}{(a+b)^2}=\frac{a}{a+b}
\end{align*}
注意这里的 $P(B\mid A)=P(B)$, 即事件 $A$ 发生与否, 对事件 $B$ 发生的概率没有影响. 从直观上讲, 这很自然. 因为我们这里采用的是有放回摸球, 因此第二次摸球时袋中球的组成与第一次摸球时完全相同, 当然第一次摸球的结果实际上不影响第二次摸球. 在这种场合可以说, 事件 $A$ 与事件 $B$ 的出现有某种"独立性".
\end{proof}
\end{example}

对此, 我们引进

\begin{definition}[事件的独立性]\label{definition:ch02-事件独立性}
对事件 $A$ 及 $B$, 若
\begin{align}\label{align:ch02-两事件独立定义}
P(AB)=P(A)P(B)
\end{align}
则称它们是统计独立的, 简称独立的（independent）.
\end{definition}

注意, 按照这个定义, 必然事件 $\Omega$ 及不可能事件 $\varnothing$ 与任何事件独立. 此外, 从\eqref{align:ch02-两事件独立定义}中看出, $A$ 与 $B$ 的位置对称, 因此亦称 $A$ 与 $B$ 相互独立.

\begin{corollary}[独立事件的条件概率]\label{corollary:ch02-独立事件条件概率}
若事件 $A$, $B$ 独立, 且 $P(B)>0$, 则
\begin{align}\label{align:ch02-独立事件条件概率}
P(A\mid B)=P(A)
\end{align}

\begin{proof}
由条件概率定义及\eqref{align:ch02-两事件独立定义}得
\begin{align*}
P(A\mid B)=\frac{P(AB)}{P(B)}=\frac{P(A)P(B)}{P(B)}=P(A)
\end{align*}
因此, 若事件 $A$, $B$ 相互独立, 则 $A$ 关于 $B$ 的条件概率等于无条件概率 $P(A)$, 这表示 $B$ 的发生对于事件 $A$ 是否发生没有提供任何信息, 独立性就是把这种关系从数学上加以严格定义.
\end{proof}
\end{corollary}

\begin{corollary}[独立事件的补运算]\label{corollary:ch02-独立事件补运算}
若事件 $A$ 与 $B$ 独立, 则下列各对事件也相互独立：
\begin{align*}
\{\bar{A},B\},\ \{A,\bar{B}\},\ \{\bar{A},\bar{B}\}
\end{align*}

\begin{proof}
由于
\begin{align*}
P(\bar{A}B)&=P(B-AB)=P(B)-P(AB)\\
&=P(B)-P(A)P(B)=P(B)[1-P(A)]\\
&=P(\bar{A})P(B)
\end{align*}
所以 $\bar{A}$ 与 $B$ 相互独立, 由它立刻推出 $\bar{A}$ 与 $\bar{B}$ 相互独立, 由 $\bar{\bar{A}}=A$ 又推出 $A$, $\bar{B}$ 相互独立.
\end{proof}
\end{corollary}

不放回摸球模型提供了不独立的一个简单的例子.

\begin{example}[不放回摸球]\label{example:ch02-不放回摸球}
在前例中, 若采用不放回摸球, 试求同样那两个事件的概率.

\begin{proof}[解]
这时
\begin{align*}
P(A)=\frac{a}{a+b},\quad P(AB)=\frac{a(a-1)}{(a+b)(a+b-1)}
\end{align*}
\begin{align*}
P(\bar{A}B)=\frac{ba}{(a+b)(a+b-1)}
\end{align*}
所以
\begin{align*}
P(B\mid A)=\frac{P(AB)}{P(A)}=\frac{a-1}{a+b-1}
\end{align*}
而
\begin{align*}
P(B)=P(AB)+P(\bar{A}B)=\frac{a}{a+b}
\end{align*}
这里 $P(B\mid A)\neq P(B)$, 即事件 $B$ 与事件 $A$ 不是相互独立的. 因为第一次摸得黑球, 事实上已使袋中球的组成成分改变了, 当然要影响第二次摸得黑球的概率. 而 $P(A)=P(B)$, 即抽签与顺序无关.
\end{proof}
\end{example}

\subsection{多个事件的独立性}

我们先定义三个事件 $A$, $B$, $C$ 的独立性.

\begin{definition}[三个事件的独立性]\label{definition:ch02-三事件独立性}
对于三个事件 $A$, $B$, $C$, 若下列四个等式同时成立, 则称它们相互独立.
\begin{subequations}\label{align:ch02-三事件独立}
\begin{align}
P(AB)&=P(A)P(B)\\
P(BC)&=P(B)P(C)\\
P(AC)&=P(A)P(C)\\
P(ABC)&=P(A)P(B)P(C)
\end{align}
\end{subequations}
按两个事件独立性的定义, 我们知道若\eqref{align:ch02-三事件独立}的前三式成立, 则 $A$ 与 $B$, $B$ 与 $C$, $C$ 与 $A$ 都相互独立, 即 $A$, $B$, $C$ 两两独立.
\end{definition}

读者自然会提出这样一个问题：三个事件 $A$, $B$, $C$ 两两独立, 能否保证它们相互独立呢？即能否由前三式推出最后一式？回答是否定的, 这从下面简单的例子就可看出.

\begin{example}[Burnstone 反例：均匀正四面体]\label{example:ch02-四面体反例}
一个均匀的正四面体, 其第一面染成红色, 第二面染成白色, 第三面染成黑色, 而第四面同时染上红, 白, 黑三种颜色. 现在以 $A$, $B$, $C$ 分别记投一次四面体出现红, 白, 黑颜色朝下的事件, 则由于在四面体中有两面有红色, 因此
\begin{align*}
P(A)=\frac{1}{2}
\end{align*}
同理 $P(B)=P(C)=\dfrac{1}{2}$, 容易算出
\begin{align*}
P(AB)=P(BC)=P(AC)=\frac{1}{4}
\end{align*}
所以前三式成立, 即 $A$, $B$, $C$ 两两独立, 但是
\begin{align*}
P(ABC)=\frac{1}{4}\neq\frac{1}{8}=P(A)P(B)P(C)
\end{align*}
因此最后一式不成立, 从而 $A$, $B$, $C$ 不相互独立.
\end{example}

下面再提供一个例子说明由最后一式不能推出前三式, 进一步说明要 $A$, $B$, $C$ 相互独立必须同时要求全部四个等式.

\begin{example}[均匀正八面体]\label{example:ch02-正八面体反例}
若有一个均匀正八面体, 其第 $1,2,3,4$ 面染红色, 第 $1,2,3,5$ 面染白色, 第 $1,6,7,8$ 面染上黑色, 现在以 $A$, $B$, $C$ 分别表示投一次正八面体出现红, 白, 黑的事件, 则
\begin{align*}
P(A)=P(B)=P(C)=\frac{4}{8}=\frac{1}{2}
\end{align*}
\begin{align*}
P(ABC)=\frac{1}{8}=P(A)P(B)P(C)
\end{align*}
但是
\begin{align*}
P(AB)=\frac{3}{8}\neq\frac{1}{4}=P(A)P(B)
\end{align*}
\end{example}

现在我们可以定义 $n$ 个事件的独立性.

\begin{definition}[$n$ 个事件的独立性]\label{definition:ch02-n事件独立性}
对 $n$ 个事件 $A_1, A_2, \cdots, A_n$, 若对于所有可能的组合 $1\leqslant i<j<k<\cdots\leqslant n$ 成立着
\begin{subequations}\label{align:ch02-n事件独立}
\begin{align}
P(A_iA_j)&=P(A_i)P(A_j)\\
P(A_iA_jA_k)&=P(A_i)P(A_j)P(A_k)\\
&\cdots\cdots\nonumber\\
P(A_1A_2\cdots A_n)&=P(A_1)P(A_2)\cdots P(A_n)
\end{align}
\end{subequations}
则称 $A_1, A_2, \cdots, A_n$ 相互独立.
\end{definition}

这里第一行有 $\binom{n}{2}$ 个式子, 第二行有 $\binom{n}{3}$ 个式子等等, 因此共应满足
\begin{align*}
\binom{n}{2}+\binom{n}{3}+\cdots+\binom{n}{n}=2^n-n-1
\end{align*}
个等式. 由三个事件的场合可看出同时满足这些关系式是必须的.

显然若 $n$ 个事件相互独立, 则它们中的任何 $m$（$2\leqslant m<n$）个事件也相互独立. 此外对于多个相互独立的事件也成立着类似于上述推论 1 及推论 2 的结果, 读者试自行叙述并验证之.

最后, 称无穷多个事件是相互独立的, 如果其中任意有限多个事件都相互独立.

\subsection{事件独立性与概率的计算}

从事件独立性的定义立刻能看出：若事件是独立的, 则许多概率的计算就可以大为简化. 下面先举两个例子.

\begin{remark}[相互独立事件至少发生其一的概率的计算]\label{remark:ch02-独立事件至少发生其一}
若 $A_1, A_2, \cdots, A_n$ 是 $n$ 个相互独立的事件, 则由于
\begin{align*}
\overline{A_1\cup A_2\cup\cdots\cup A_n}=\bar{A}_1\bar{A}_2\cdots\bar{A}_n
\end{align*}
因此
\begin{align}\label{align:ch02-独立事件至少发生其一}
P(A_1\cup A_2\cup\cdots\cup A_n)&=1-P(\bar{A}_1\bar{A}_2\cdots\bar{A}_n)\\
&=1-P(\bar{A}_1)P(\bar{A}_2)\cdots P(\bar{A}_n)\nonumber
\end{align}
这个公式比起不独立的场合, 要简便得多, 它经常被用到.
\end{remark}

\begin{example}[混合血清中的肝炎病毒]\label{example:ch02-混合血清肝炎病毒}
假若每个人血清中含有肝炎病毒的概率为 $0.4\%$, 混合 $100$ 个人的血清, 求此血清中含有肝炎病毒的概率.

\begin{proof}[解]
以 $A_i$（$i=1,2,\cdots,100$）记第 $i$ 个人的血清含有肝炎病毒这一事件, 可以认为它们相互独立, 所求概率为 $P(A_1\cup\cdots\cup A_{100})$, 由\eqref{align:ch02-独立事件至少发生其一}得
\begin{align*}
P(A_1\cup\cdots\cup A_{100})=1-P(\bar{A}_1)\cdots P(\bar{A}_{100})=1-0.996^{100}\approx 0.33
\end{align*}
虽然每个人有病毒的概率很小, 但是混合后则有很大的概率, 在实际工作中, 这类效应值得充分重视.
\end{proof}

顺带指出, 没有独立性的假定, 上述计算便无从进行. 当然这里的独立性只能是一种近似. 当把某种数学模型用于实际问题时, 这种近似是不可避免的. 因此, 作理论研讨时, 独立性必须按定义验证; 解决实际问题时, 独立性通常只是一种恰当的假定.
\end{example}

\begin{remark}[在可靠性理论中的应用]\label{remark:ch02-可靠性理论应用}
对于一个元件, 它能正常工作的概率 $p$, 称为它的可靠性. 元件组成系统, 系统正常工作的概率称为该系统的可靠性. 随着近代电子技术的迅猛发展, 关于元件和系统可靠性的研究已发展成为一门新的学科——可靠性理论.
\end{remark}

这里, 我们通过一些例子来说明有关的概念.

\begin{example}[附加通路系统的可靠性]\label{example:ch02-附加通路可靠性}
如果构成系统的每个元件的可靠性均为 $r$, $0<r<1$, 且各元件能否正常工作是相互独立的, 试求下面附加通路系统的可靠性（见图\ref{figure:ch02-附加备份系统}(a)）.

\begin{proof}[解]
每条通路要能正常工作, 当且仅当该通路上各元件正常工作, 故其可靠性为
\begin{align*}
R_c=r^n
\end{align*}
即通路发生故障的概率为 $1-r^n$. 由于系统是由两通路并联而成的, 两通路同时发生故障的概率为 $(1-r^n)^2$, 因此附加通路系统的可靠性为
\begin{align*}
R_s=1-(1-r^n)^2=r^n(2-r^n)=R_c(2-R_c)
\end{align*}
注意到 $R_c<1$, 故 $R_s>R_c$, 所以附加通路能使系统的可靠性增加.
\end{proof}
\end{example}

\begin{figure}[H]
\centering
\includegraphics[width=0.5\textwidth]{ch02-fig2-2-1.png}
\caption{附加备份系统}
\label{figure:ch02-附加备份系统}
\end{figure}

\begin{example}[附加元件系统的可靠性]\label{example:ch02-附加元件可靠性}
在前例条件下, 试求下面附加元件系统的可靠性（见图\ref{figure:ch02-附加备份系统}(b)）.

\begin{proof}[解]
每对并联元件的可靠性为
\begin{align*}
R'=1-(1-r)^2=r(2-r)
\end{align*}
系统由各对并联元件串联而成, 故其可靠性为
\begin{align*}
R_s'=(R')^n=r^n(2-r)^n
\end{align*}
显然 $R_s'>R_c$. 因此用附加元件的方法同样也能增加系统的可靠性.

利用数学归纳法不难证明当 $n\geqslant 2$ 时, $(2-r)^n>2-r^n$, 即 $R_s'>R_s$. 因此虽然上面两个系统同样由 $2n$ 个元件构成, 作用也相同, 但是第二种构成方式比第一种方式可靠性来得大. 寻找可靠性较大的构成方式也是可靠性理论的研究课题之一.
\end{proof}

从上述讨论可以看出, 元件与系统的可靠性是用概率来定义的, 所以概率论是研究可靠性理论的重要工具.
\end{example}

\subsection{试验的独立性}

有了事件独立性的概念, 我们可以定义试验的独立性. 直观上, 试验 $E_1$ 与试验 $E_2$ 要能说是独立的, 应是指试验 $E_1$ 的结果的发生与试验 $E_2$ 的结果的发生是独立的, 所以自然想到通过各试验的事件间的独立性来定义试验的独立性. 为了做到这点, 首先要构造一个能描述这些试验的公共的样本空间.

设试验 $E_1$ 的样本空间是 $\Omega_1=\{\omega^{(1)}\}$, 试验 $E_2$ 的样本空间是 $\Omega_2=\{\omega^{(2)}\}$, $\cdots\cdots$ 试验 $E_n$ 的样本空间是 $\Omega_n=\{\omega^{(n)}\}$, 为描述这 $n$ 次试验, 应构造复合试验 $E$, 它表示依次进行试验 $E_1, E_2, \cdots, E_n$, 其样本点为
\begin{align*}
\omega=(\omega^{(1)},\omega^{(2)},\cdots,\omega^{(n)})
\end{align*}
这个样本空间为 $\Omega_1, \Omega_2, \cdots, \Omega_n$ 的乘积空间, 记作
\begin{align*}
\Omega=\Omega_1\times\Omega_2\times\cdots\times\Omega_n
\end{align*}

\begin{example}[掷硬币与摸球的复合试验]\label{example:ch02-复合试验}
若试验 $E_1$ 是掷一枚硬币, $\Omega_1=\{$正, 反$\}$, 试验 $E_2$ 是从装有红白黑三球的袋子中摸出一球, $\Omega_2=\{$红, 白, 黑$\}$, 则复合试验 $E$ 表示先掷一枚硬币再摸一球, 它相应的样本空间 $\Omega=\Omega_1\times\Omega_2$ 由下列 $6$ 个样本点构成：（正, 红）, （正, 白）, （正, 黑）, （反, 红）, （反, 白）, （反, 黑）.
\end{example}

接着我们可以引进"与第 $k$ 次试验有关的事件"的概念, 这种事件发生与否仅与第 $k$ 次试验的结果有关, 即为了判断某一样本点是否属于这个事件, 只需察看它的第 $k$ 个分量. 例如在例\ref{example:ch02-复合试验}中, "第 $2$ 次试验摸出红球"这一事件就是一个与第 $2$ 次试验有关的事件. 值得指出, 必然事件 $\Omega$ 及不可能事件 $\varnothing$ 可以认为与所有的试验有关.

现在若以 $\mathscr{A}_k$ 记与第 $k$ 次试验有关的事件全体, 则可以通过下列方式定义试验的独立性.

\begin{definition}[试验的独立性]\label{definition:ch02-试验独立性}
若对于任意的
\begin{align*}
A^{(1)}\in\mathscr{A}_1,\ A^{(2)}\in\mathscr{A}_2,\ \cdots,\ A^{(n)}\in\mathscr{A}_n
\end{align*}
均成立
\begin{align*}
P(A^{(1)}A^{(2)}\cdots A^{(n)})=P(A^{(1)})P(A^{(2)})\cdots P(A^{(n)})
\end{align*}
则称试验 $E_1, E_2, \cdots, E_n$ 是相互独立的.
\end{definition}

注意到 $\Omega\in\mathscr{A}_i$, $i=1,2,\cdots,n$, 因此由定义立刻推出, 若 $n$ 个试验相互独立, 则其中的 $m$（$2\leqslant m<n$）个试验也是相互独立的.

在例\ref{example:ch02-复合试验}中, 若对样本空间 $\Omega$ 中的 $6$ 个样本点都给定概率 $\dfrac{1}{6}$, 则不难验证试验 $E_1$ 与试验 $E_2$ 是相互独立的. 读者可以思考一下, 是否还可以给定其他概率使试验 $E_1$ 与试验 $E_2$ 独立？

$n$ 次有放回摸球所构成的 $n$ 个试验是相互独立的, 而 $n$ 次不放回摸球模型则是 $n$ 个试验不独立的简单例子.

特别重要的一类试验是所谓重复独立试验, 这时 $\Omega_1=\Omega_2=\cdots=\Omega_n$, 有关事件的概率保持不变, 而且各次试验是相互独立的. 投 $n$ 次硬币或进行 $n$ 次有放回摸球是重复独立试验的简单例子. 重复独立试验是作为"在同样条件下重复试验"的数学模型而出现的, 它在概率论中很有地位, 因为随机现象的统计规律性只有在大量重复试验中才会显示出来.

我们在下一节将研究一类最简单的重复独立试验——Bernoulli（J. Bernoulli, 1654—1705）试验.

\end{document}